\section{Background and Related Work}\label{sec:background}

\subsection{Matrix--Vector Products in a ReRAM Crossbar}
A crossbar places a resistive cell at every crossing of a row (word line) and a column (bit line). If
cell $(i,j)$ has conductance $G_{ij}$ and row $i$ is driven at voltage $V_i$, the current collected at
column $j$ is
\begin{equation}
I_j=\sum_i V_i\,G_{ij},
\label{eq:crossbar}
\end{equation}
by Ohm's law in each cell and Kirchhoff's current law on the column. One read therefore computes a full
matrix--vector product in the analog domain. Accelerators such as ISAAC~\cite{shafiee2016isaac} and
PRIME~\cite{chi2016prime} build on this primitive. Three conventions from these designs are used
throughout this report:
\begin{itemize}
\item Conductances are non-negative, so a signed weight is stored as the difference of two cells, a
\emph{differential pair}.
\item A cell holds only a few bits, so a multi-bit weight is split into \emph{bit slices} on separate
cells; the column results of the slices are shifted and added digitally.
\item Arrays have a fixed size (we use $128\times128$ tiles), and an analog-to-digital converter (ADC)
digitises each column current.
\end{itemize}

\subsection{SLC, MLC and Non-Idealities}
A cell programmed to $L=2^b$ levels between $G_{\min}$ and $G_{\max}$ stores $b$ bits; $b=1$ is a
single-level cell (SLC) and $b\geq2$ a multi-level cell (MLC). The level spacing is
$\Delta=(G_{\max}-G_{\min})/(L-1)$, so a two-bit cell has levels three times closer than an SLC. A
conductance error of fixed size is therefore three times larger relative to an MLC level, while MLC
needs fewer cells per weight. Real arrays deviate from the programmed state in several ways:
\begin{itemize}
\item \emph{Programming variation}: each cell lands near, not on, its target conductance; the error is
fixed once the chip is programmed.
\item \emph{Read noise}: the current fluctuates at every read.
\item \emph{Stuck-at faults}: some cells are permanently stuck at $G_{\max}$ or $G_{\min}$, whatever
is written; faults can be scattered, clustered, or affect a whole line.
\item \emph{IR drop}: wire resistance lowers the voltage reaching cells far from the drivers.
\item \emph{ADC resolution}: column currents are rounded to a limited number of levels.
\end{itemize}

\subsection{HyFlexPIM}
HyFlexPIM~\cite{song2025hyflexpim} is a mixed-signal processing-in-memory accelerator for
transformers. Its analog part can switch each array between SLC and MLC operation, so that critical
weights use SLC and the rest use denser, cheaper MLC. Placing raw weights this way gives little benefit,
because importance is spread over many weights. HyFlexPIM therefore factorises each weight matrix with a
truncated SVD, $W\approx U\Sigma V^{\top}$, and fine-tunes the factors, which concentrates the loss
gradient in a small subset of components; it calls this \emph{gradient redistribution}. Components are
ranked by the gradient of the loss with respect to their singular value, and the top-ranked share is
mapped to SLC. Accuracy is evaluated with relative Gaussian noise applied to whole weights,
$\tilde W=W\odot(1+\eta)$, at a single level derived from a measured MLC bit error
rate~\cite{fan2024mlc}, together with a limited ADC resolution. Stuck cells, read noise and IR drop are
not modelled. Our work keeps the SVD
factorisation and the hybrid layout unchanged; \cref{sec:intro} lists what it adds.

\subsection{Fault-Tolerant Mapping and Training}
Two families of methods make crossbar inference tolerant to device errors. The first adapts training to
the hardware. Chen \emph{et al.}~\cite{chen2017accelerator} include device variations and defects in
training. Xia \emph{et al.}~\cite{xia2017fault} detect stuck cells on line and adapt training to them.
Rasch \emph{et al.}~\cite{rasch2023hwa} inject hardware noise during training at scale, in what is now
called hardware-aware training; gradients pass through the non-differentiable hardware model with a
straight-through estimator~\cite{bengio2013ste}. The second family leaves the weights unchanged and
chooses \emph{where} they are stored, so that important values avoid faulty cells. Such mapping
problems reduce to assignment problems, which the Hungarian method~\cite{kuhn1955hungarian} solves
exactly.

Our fault-aware mapping belongs to the second family, with two differences. First, the freedom to
reorder comes from the SVD itself: the components of a factorised layer can be stored in any order
without changing its output, so no permutation hardware or change to neighbouring layers is needed.
Second, the cost of a placement is weighted by the Fisher information of each stored value, a
second-order estimate of loss sensitivity in the spirit of Optimal Brain
Damage~\cite{lecun1990obd} and of Fisher-based importance in continual
learning~\cite{kirkpatrick2017ewc}.
